\setcounter{chapter}{-1}

\chapter{绪论}
\label{ch:intro}

把氢原子和氦原子发出的光分别送入光谱仪，会看到两套不同的亮线。天文学家正是凭各元素特有的谱线，从星光中辨认恒星大气的成分。是什么决定了一条谱线的波长？
量子力学(quantum mechanics)研究物质和辐射在原子及亚原子尺度上遵循的规律。原子能发出什么频率的光，几个原子能否结合成稳定的分子，都是它要回答的问题。

按经典力学计算，电子可以沿不同半径的圆轨道绕核运动，轨道能量随半径连续变化。氢原子的束缚态却只有一系列分立的能量，两个能级之差决定了原子在它们之间跃迁时发出或吸收的光的频率。这就是能量量子化的一个例子。脱离原子后，自由电子的动能仍可连续取值。

经典力学还告诉我们，给定作用力和初始的位置、动量，质点的轨迹就随之确定。若把电子制备在相同的量子态，逐个送过双缝，电子却会落在屏上的不同位置。
量子力学预言这些落点的概率分布。两条缝同时开放时，经两条路径到达同一处的概率振幅可以相互抵消，屏上便出现暗纹。在这些位置，打开第二条缝反而降低了探测到电子的概率。

这些微观规律也支配着日常所见的材料和器件。量子化学从电子态算出拆开一个化学键需要多少能量。半导体的能带结构决定发光二极管的发光颜色，也是理解微处理器中晶体管如何控制电流的基础。
超导磁体中，电子的集体运动使电流能够长期维持，产生稳定的强磁场。
量子光学研究光与物质的相互作用，从激光器里的受激辐射到单个光子的制备和探测。一个光子的偏振就可以编码一个量子比特。量子信息科学研究如何用这样的量子态传输和处理信息，量子通信和量子计算都是这方面的量子技术。
把量子原理用于场，还能描述粒子的产生和湮灭。在宇宙学中，人们用量子场论研究早期宇宙的量子涨落。

\section{量子力学发展简史}
\label{sec:timeline}
量子力学并不只是 \textit{shut up and calculate} 的计算工具。这套理论既解释了已有的实验结果，也预言了当时尚未观察到的现象。

十九世纪末，原子物理学家们已在不同元素的发射光谱中辨认出各自的分立谱线(discrete spectral lines)。
Balmer 用一个含整数的经验公式概括了氢的几条可见谱线的波长。
经典的周期运动也能产生分立的辐射频率，但经典原子模型允许电子在不同半径上以不同频率绕核运动，无法选出氢特有的这组频率。
原子的稳定性也是个问题。按经典电动力学，绕核运动的电子会辐射能量，轨道随之收缩。

一条线索来自光。1900 年，Planck 为解释黑体辐射谱引入了能量元。1905 年，Einstein 提出光量子假说，由此预言光电子的最大动能随入射光频率线性增加（\ssecref{sec:photons}）。
1913 年，Bohr 把光的频率与原子的能量联系起来。他假定电子可以在某些轨道上稳定运动，原子从高能轨道跃迁到低能轨道时，轨道间的能量差决定发光频率。他由此算出了 Balmer 线系的波长\cite{bohrlecture}。
但 Bohr 模型无法解释氢的精细结构，也无法算出氦原子和更复杂原子的谱线\cite{helium-models}。

量子力学理论就是为了解释这些实验现象。
1925--1926 年，矩阵力学和波动力学相继建立。求氢原子的能级时，只需求解 Schrödinger 方程。哪些束缚态可以存在、各自的能量是多少，都由方程和波函数的边界条件决定（\ssecref{ss:qm-old-quantum}）。
换成氦原子，还须计入两个电子之间的排斥。1929 年，Hylleraas 用两电子波函数做近似计算，得到的电离能已与实验值很接近\cite{helium-models}。

量子力学理论还能预言尚待观察的现象。1923--1924 年，de Broglie 提出动量大小为 $p$ 的电子具有波长 $h/p$，其中 $h$ 是 Planck 常数\cite{debroglie}。
如果这个波长有物理意义，用电子束照射晶面间距已知的晶体，就能预先算出哪些方向会出现衍射峰。1927 年，Davisson 和 Germer 在镍单晶上测出的峰位与这个波长相符（\secref{sec:interference}）\cite{davisson}。
按照 Born 在 1926 年给出的概率解释，电子波在某个方向的强度越大，该方向探测到电子的概率就越大。1928 年，Dirac 又把狭义相对论纳入电子的波动方程（\chapref{ch:relativity}）。

\tabref{tab:history}汇集了本卷讨论的实验，按年代排序。
节号指向相应章节，实验名和页码可直接跳到讲解该实验的段落。

\begin{notecaptionblock}
    \captionof{table}{量子力学实验与观测索引}\label{tab:history}
    \small
    \setlength{\tabcolsep}{1.4pt}
    \renewcommand{\arraystretch}{1.2}
    \newcommand{\experimentref}[3]{\hyperref[#1]{#3} & \ref{#2} & \pageref{#1}}
    \begin{tabular*}{\linewidth}{@{\extracolsep{\fill}}l>{\raggedright\arraybackslash}p{.25\linewidth}rr@{\hspace{1.2em}}l>{\raggedright\arraybackslash}p{.28\linewidth}rr@{}}
        \toprule
        年份 & 实验与观测 & 节号 & 页码 & 年份 & 实验与观测 & 节号 & 页码\\
        \midrule
        1885 & \experimentref{exp:qm-balmer}{sec:qm-rydberg-lines}{氢的 Balmer 线系}
            & 1982 & \experimentref{exp:qm-bell-1982}{ss:chsh-derivation}{快速切换分析器的 Bell 检验}\\
        1887 & \experimentref{sec:photons}{sec:photons}{Hertz 紫外光放电实验}
            & 1985 & \experimentref{exp:qm-junction-escape}{sec:qm-macroscopic}{超导结量子隧穿}\\
        1896 & \experimentref{ss:zeeman-experiment}{ss:zeeman-experiment}{Zeeman 谱线展宽}
            & 1985 & \experimentref{exp:qm-junction-levels}{sec:qm-macroscopic}{超导结能级的微波测量}\\
        1899--1900 & \experimentref{ss:blackbody-experiment}{ss:blackbody-experiment}{黑体长波辐射测量}
            & 1986 & \experimentref{exp:qm-photon-anticorrelation}{sec:photoncount}{单光子分束反关联}\\
        1914 & \experimentref{sec:fh}{sec:fh}{Franck--Hertz 汞蒸气实验}
            & 1987 & \experimentref{ss:berry-readout}{ss:berry-readout}{中子自旋的几何相位}\\
        1916 & \experimentref{exp:qm-millikan-stopping}{sec:photons}{Millikan 光电截止电压}
            & 1987 & \experimentref{sec:hom}{sec:hom}{Hong--Ou--Mandel 干涉}\\
        1922 & \experimentref{ss:sg-apparatus}{ss:sg-apparatus}{Stern--Gerlach 银原子束}
            & 1989 & \experimentref{exp:qm-tonomura}{sec:interference}{外村彰单电子干涉实验}\\
        1923 & \experimentref{ss:compton}{ss:compton}{Compton 散射}
            & 1995 & \experimentref{exp:qm-absorber-1995}{ss:qm-absorber-test}{单光子干涉仪检出吸收体}\\
        1927 & \experimentref{exp:qm-nickel-diffraction}{sec:interference}{镍单晶电子衍射}
            & 1998 & \experimentref{exp:qm-path-1998}{ss:qm-path-marking}{原子干涉仪路径标记}\\
        1927 & \experimentref{exp:qm-film-diffraction}{sec:interference}{薄膜电子衍射}
            & 2002 & \experimentref{exp:qm-eraser-2002}{ss:qm-quantum-eraser}{双缝量子擦除实验}\\
        1938 & \experimentref{ss:rabi-measurement}{ss:rabi-measurement}{Rabi 分子束磁共振\textsuperscript{a}}
            & 2012 & \experimentref{exp:qm-neutron-disturbance}{ss:qm-error-disturbance}{中子自旋的顺序测量}\\
        1947 & \experimentref{ss:lamb-experiment}{ss:lamb-experiment}{Lamb--Retherford 氢谱}
            & 2015 & \experimentref{exp:qm-bell-2015}{ss:chsh-derivation}{远距电子自旋的 Bell 检验}\\
        1960 & \experimentref{sec:ab}{sec:ab}{Chambers 磁通干涉实验}
            & 2023 & \experimentref{exp:qm-bild-cat}{ss:qm-schrodinger-cat}{机械振子的猫态}\\
        1972 & \experimentref{exp:qm-bell-1972}{ss:chsh-derivation}{光子偏振关联的 Bell 检验}
            & 2024 & \experimentref{exp:qm-surface-memory}{sec:qm-error-correction}{表面码的逻辑存储}\\
        1975 & \experimentref{exp:qm-cow-gravity}{sec:qm-gravity}{中子引力干涉}
            & 2025 & \experimentref{exp:qm-recoil-2025}{ss:qm-path-marking}{单原子反冲与光子干涉}\\
        1980 & \experimentref{sec:qm-integer-hall}{sec:qm-integer-hall}{整数量子 Hall 效应}
            & 2026 & \experimentref{exp:qm-cluster-interference}{sec:qm-macroscopic}{钠团簇的光栅干涉}\\
        \bottomrule
    \end{tabular*}\par
    \smallskip
    {\footnotesize\raggedright
        \textsuperscript{a}\,Ramsey 于 1949 年提出让分子束先后通过两个分开的振荡场区域，1950 年发表了这一方法的系统论述（\secref{sec:ramsey}）\cite{ramsey}。\par}
\end{notecaptionblock}

\section{量子理论的实验基础}
光照在金属上，或电子撞上原子时，怎样从仪器读数知道它们交换了多少能量？
光电实验用电压测出逸出电子的最大动能，Compton 实验用散射波长检验光子传给电子的动量，
Franck--Hertz 实验则让电子碰撞原子，从电流的起伏读出原子的激发能。
三种装置都可以先用大学物理的守恒律来分析，也由此看出哪些微观规律还要由实验给定。

\subsection{光电效应}
\label{sec:photons}

1887 年，Hertz 发现，用紫外光照射火花间隙的电极，原本不足以引起放电的条件下也能打出火花。后来人们才认出，紫外光可以使电子逸出金属表面\cite{history}。怎样测出逸出电子的数目和能量？
把金属阴极置于真空中，用频率为$\nu$的近单色光照射，
另一电极收集逸出的电子，电流表测量单位时间到达的电荷。
改变两电极之间的电压，可以分别测量发射率和电子的最大动能。

先让收集极吸引电子。电压足够大后，能逸出的电子几乎都被收集，
电流接近饱和值；在表面状态不变的弱光范围，饱和电流随光强增加。
再使收集极排斥电子，电子必须克服电势能障碍才能到达它。
反向电压增大到某个值时，电流降至背景，说明能量最高的电子也无法通过。
校正两电极间的接触电势后，将刚好阻止光电子到达收集极的有效反向电压大小记为截止电压(stopping potential) $V_{\stopping}$。
电荷为$-e$的电子克服的能量障碍是$eV_{\stopping}$，因而
$K_{\max} = eV_{\stopping}$，其中$e > 0$为元电荷。

在同一块金属表面，固定频率而增大光强，饱和电流增大，截止电压基本不变。改变入射频率，低于某个阈频时，单个光子无法使电子逸出。
1916 年，\phantomsection\label{exp:qm-millikan-stopping}Millikan 报告了不同入射频率下的截止电压测量。超过阈频后，截止电压随频率线性增加，直线的斜率给出 $h/e$\cite{history}。
若一次吸收向电子转移能量$h\nu$，其中$h$为Planck常数，
电子从材料逸出还要付出逸出功(work function) $W$，
则能量守恒给出光电效应(photoelectric effect)的基本关系
\begin{equation}
    eV_{\stopping} = K_{\max} = h\nu - W,\qquad e > 0.
\end{equation}
更换金属改变了电子逸出表面所需的能量 $W$；光子的能量仍由频率唯一确定（图\ref{fig:photoelectric-threshold}）。

\begin{figure}[H]
    \centering
    \raisebox{-.5\height}{\tikzfigure{photoelectric-cell}}\hfill
    \raisebox{-.5\height}{\includegraphics[width=.52\linewidth]{figures/photoelectric_threshold.pdf}}
    \caption{光电截止电压。左：波浪线为入射光，直箭头为从阴极 K 飞向收集极 C 的电子。右：各频率下使电流表 A 的读数降至背景的反向电压已经过接触电势校正，直线的斜率为 $h/e$，横轴截距为 $W/h$。}
    \label{fig:photoelectric-threshold}
\end{figure}

设入射光功率为 $\mathcal P$。每个光子的能量是 $h\nu$，每秒入射的光子数就是 $\mathcal P/(h\nu)$。
以 $\eta$ 表示吸收、逸出和收集的总效率，饱和电流近似为
\begin{equation}
    I_\mathrm{sat} = e\eta\frac{\mathcal P}{h\nu}.
\end{equation}
截止电压读出单个光电子的最大动能，饱和电流除以 $e$ 给出每秒收集的电子数。

\begin{example}{光电截止电压与入射频率}{qm-photoelectric-fit}
    在同一块稳定的金属表面，用两个不同的入射频率 $\nu_1,\nu_2$ 测得截止电压 $V_1,V_2$。怎样由读数求出 $h/e$？两次测量中相同的接触电势对结果有何影响？

    \par\smallskip\textbf{解}\quad 两个光电能量方程相减，逸出功消去，得到
    \begin{equation}
        e(V_2 - V_1) = h(\nu_2 - \nu_1).
    \end{equation}
    由两次读数的差便可求出斜率 $h/e$。
    共同的接触电势也只改变截距；若表面在两次测量之间变化，
    两个方程中的$W$便不再相同，需要重新控制表面与电极条件。
\end{example}

上述解释适用于弱光下单光子吸收占主导的范围。
强场中可以吸收多个光子，复杂表面的功函数分布也会展宽截止区。
截止电压随频率的变化给出了单次吸收的能量规律 $eV_{\stopping} + W = h\nu$；
\secref{sec:photoncount}再用符合计数区分单光子态与弱相干光。

\subsection{Compton散射}
\label{ss:compton}

1923 年，Compton 用近单色 X 射线照射石墨，再用晶体光谱仪测量散射光。除了原波长的成分，他还测到波长较长的成分。散射角越大，两种波长相差越多\cite{compton}。
散射角与波长位移的关系同时受能量守恒和动量守恒约束。
采用光子能量$E = h\nu$、动量大小$p = E/c$，
并把参与碰撞的电子在碰撞时间内近似为自由粒子，就能定量计算位移。

对初始静止的自由电子，入射、出射光子的能量分别为$E,E'$，夹角为$\theta$。
$m_\electron$是电子静质量，$c$是真空光速。
由能量守恒，反冲电子的总能量为 $E_\electron = m_\electron c^2 + E - E'$。它的动量等于光子碰撞前后的动量之差。把这个动量的大小 $p_\electron$ 代入相对论能量关系 $E_\electron^2 = m_\electron^2c^4 + p_\electron^2c^2$，得到
\begin{equation}
    \begin{alignedat}{1}
        (m_\electron c^2 + E - E')^2 &= m_\electron^2c^4 + E^2 + E'^2 - 2EE'\cos\theta,\\
        m_\electron c^2(E - E') &= EE'(1 - \cos\theta).
    \end{alignedat}
\end{equation}
展开第一行左侧的平方，
\[
    (m_\electron c^2 + E - E')^2
    = m_\electron^2c^4 + 2m_\electron c^2(E - E') + E^2 + E'^2 - 2EE' .
\]
与右侧相同的静质量项和两个平方项消去，再除以 2 便得第二行。
将第二行除以$EE'm_\electron c^2$，得到
\begin{equation}
    \frac1{E'} - \frac1E = \frac{1 - \cos\theta}{m_\electron c^2}.
\end{equation}
代入$E = hc/\lambda$、$E' = hc/\lambda'$便得
\begin{equation}
    \lambda' - \lambda = \frac{h}{m_\electron c}(1 - \cos\theta).
\end{equation}
对原子中的电子，若束缚能远小于光子转移的能量，且电子的初始动量远小于光子转移的动量，就可以采用这个自由静止电子的近似。
光子将部分能量交给反冲电子，出射波长随散射角增大而增长（图\ref{fig:compton-shift}）。各方向的计数强度还取决于散射相互作用。
\begin{figure}[htbp]
    \centering
    \raisebox{-.5\height}{\tikzfigure{compton-geometry}}\hfill
    \raisebox{-.5\height}{\includegraphics[width=.52\linewidth]{figures/compton_shift.pdf}}
    \caption{Compton 散射。左：蓝、橙波浪线分别标记入射和出射 X 射线，$\theta$ 从入射束的延长线量起；直箭头为反冲电子。右：横轴取 $\theta$，纵轴以电子的 Compton 波长 $\lambda_{\mathrm{C}} = h/(m_\electron c)$ 为位移单位。}
    \label{fig:compton-shift}
\end{figure}

\subsection{Franck--Hertz实验}
\label{sec:fh}
原子的分立谱线对应束缚能级之间的跃迁。若改用电子碰撞原子，电子至少需要多少动能才能激发原子？
Franck--Hertz实验用加热阴极发射电子，以可调加速电压$U$使电子穿过稀薄汞蒸气，
再由栅极后的收集极测量电流。栅极与收集极之间保留一个小的反向电压$U_{\retarding}$，
使刚刚损失了大部分动能的电子不能到达收集极\cite{franck}。

先考虑电子能量低于第一个显著激发阈值的情况。
电子与重原子的弹性碰撞主要改变方向；由于电子质量远小于原子质量，
一次弹性碰撞转移给原子质心的能量比例很小。
增大$U$通常使更多电子克服反向势垒，收集电流增加。
若电子获得的动能刚好足以激发原子，并且碰撞发生在接近栅极的位置，
它会把一份能量$\Delta E$交给原子，剩余动能不足以克服$eU_{\retarding}$，
于是收集电流出现下降。

继续增大$U$，失能后的电子可以在剩余路程重新加速，
收集电流恢复；更高电压下又可能发生两次、三次相同类型的非弹性碰撞。
每多发生一次这样的碰撞，电子就多损失一份 $\Delta E$，相邻电流谷因而相隔约 $\Delta E/e$。把第 $r$ 个谷对应的加速电压记为 $U_r^{*}$，用 $U_{\mathrm{c}}$ 表示接触电势和收集阈值引起的共同偏移，就有
\begin{subequations}
    \begin{align}
        U_r^{*} &\simeq U_{\mathrm{c}} + r\frac{\Delta E}{e},\\
        e\bigp{U_{r+1}^{*} - U_r^{*}} &\simeq \Delta E.
    \end{align}
\end{subequations}
对图\ref{fig:franck-hertz-current}中的相邻电流谷取电压差，共同的偏移 $U_{\mathrm{c}}$ 便消去。
\begin{figure}[htbp]
    \centering
    \raisebox{-.5\height}{\tikzfigure{franck-hertz-tube}}\hfill
    \raisebox{-.5\height}{\includegraphics[width=.52\linewidth]{figures/franck_hertz_current.pdf}}
    \caption{Franck--Hertz 管及收集电流。左：K、G、C 下方依次标出热阴极、栅极、收集极的电势。右图给出收集电流随加速电压的变化，产生了多个电流谷。}
    \label{fig:franck-hertz-current}
\end{figure}

汞实验中约$\qty{4.9}{V}$的重复间距对应约$\qty{4.9}{eV}$的能量损失。
若原子随后由该激发态回到基态并辐射一个光子，
能量关系预言
\begin{equation}
    \lambda = \frac{hc}{\Delta E} \simeq \qty{253}{nm},
\end{equation}
与相应的汞紫外谱线处于同一尺度。
碰撞阈值与光谱能量差由此互相核对：前者通过电子失能读出，
后者通过发射光的频率读出。

温度会改变汞蒸气压。蒸气压太低，碰撞太少，电流谷难以分辨。蒸气压太高，多次碰撞会使峰谷变得模糊。
要进一步预言峰谷的深度和形状，还要知道阴极发射电子的初速度分布、电子在汞蒸气中的碰撞截面，以及收集极的反向电压。

把碰撞前、后原子所在的能级分别记作 $E_\initial,E_\final$。同一对能级也能用光谱来比较。
忽略原子整体反冲时，吸收或发射光子的频率满足
\begin{equation}
    h\nu = \abs{E_\final - E_\initial}.
    \label{eq:spectral-frequency}
\end{equation}
能级差给出谱线的位置。哪些跃迁容易发生，还取决于光与原子的相互作用。

\section{量子力学从哪些假设出发}
Bohr 模型和 Schrödinger 理论都能给出氢原子的主要谱线，两者对基态轨道角动量的预言却不同（\ssecref{ss:qm-old-quantum}）。
还可以把氢换成氘，检验理论能否仅凭核质量的变化就预言谱线的移动。

怎样从一个态算出测量概率？量子力学用复向量表示纯态。以双缝干涉为例，一个电子通过任一条缝都能到达屏上的同一位置。若两条路径保持相干，就把两条路径到达该处的复振幅相加，再按 Born 规则取总振幅的模平方，得到该处的概率密度（假设~\ref{psl:qm-born}及\secref{sec:interference}）。
只改变一路的相位，即使两路单独开放时的计数都不变，两路同时开放时的干涉条纹也会移动。

封闭系统的演化保持叠加关系和总概率，并且可逆，这样的变换称为幺正变换。
再假定能量决定各分量积累相位的速率，就能从短时演化推出\chapref{ch:dynamics}的 Schrödinger 方程。
这些假设的物理依据来自实验。具体系统还要给出相互作用，例如计算氢的低能谱时，要用电子与核的质量、电荷和 Coulomb 势。
一旦给定这些输入，束缚态的能量和大小就由同一个方程确定。

可以用两个能级的叠加来测量相位积累的快慢。设一个原子的两个能级分别为 $E_1$ 和 $E_2$。先将原子制备在两能级等权、同相的叠加态，让它自由演化时间 $t$，再测量它是否仍处于初态。
能量为 $E_j$ 的分量积累了相位因子 $\e^{-\i E_jt/\hbar}$，其中 $\hbar = h/(2\pi)$。该分量在初态中的系数为 $1/\sqrt{2}$，投影回初态时再乘一个 $1/\sqrt{2}$。两项相加后取模平方，得到测出初态的概率
\begin{equation}\label{eq:qm-two-energy-return}
    \P_{\mathrm{return}}(t)
    = \frac14\abs{\e^{-\i E_1t/\hbar} + \e^{-\i E_2t/\hbar}}^2
    = \cos^2\biggp{\frac{(E_2 - E_1)t}{2\hbar}} .
\end{equation}
若已用光谱测出能级差，便能事先算出计数振荡的周期 $h/\abs{E_2 - E_1}$。若制备的是两种能量态各占一半的随机混合，
沿用同一台分析器，测出初始叠加态的概率便始终为 $1/2$。因此两种制备虽然给出相同的能量分布，却可以通过计数是否随时间振荡来区分。

\section{全书结构与记号约定}
\label{sec:roadmap}
\chapref{ch:spin}从自旋筛选开始。挡住一个出口，留下的原子进入第二台分析器。转动第二台分析器，两路计数的比例随之改变。
\chapref{ch:observables}用可观测量和制备态计算能量、位置等读数的概率分布。
若两台分析器同向，在它们之间加一段横向磁场，出口计数就会随原子在场中的飞行时间振荡。\chapref{ch:dynamics}用时间演化计算这个振荡周期。

\chapref{ch:oscillator}至\chapref{ch:transitions}由能级差预言谱线位置，
再由光与原子的相互作用计算哪些跃迁更强。\chapref{ch:em}至\chapref{ch:scattering}把读数换成干涉条纹、双粒子符合和不同方向的散射计数。
最后三章追问，缓慢转动磁场会积累什么相位，只观察原子会丢失什么信息，提高光谱分辨率又会看到哪些能级分裂。

在\ssecref{ss:sg-sequence}中，1000 个制备在 $+z$ 态的原子约有一半通过 $+x$ 出口，其中又约有一半到达 $-z$ 出口，最终计数的期望值为 250。
换成\secref{sec:continuous}的位置探测器，若概率密度在一个像素内近似恒定，该像素的计数比例就近似为概率密度乘以像素宽度。
\ssecref{ss:position-well}用硬墙条件确定驻波，再对探测区间积分求概率。\secref{sec:hydrogen-1s}用原点和无穷远处的条件确定 Coulomb 基态的空间尺度和能量。

在整条直线上对波包作 Fourier 变换或分部积分时，默认波函数光滑，它及所用的导数在远处充分衰减，如 Gauss 波包。带锐利边缘的反例见\ssecref{ss:qm-regularity}。

读\chapref{ch:spin}只需会用两行两列的矩阵乘列向量。自旋态的两个分量都是复数。复数 $c = u + \i v = r\e^{\i\varphi}$ 的模 $r$ 是它在复平面上到原点的距离，
相位 $\varphi$ 是从正实轴量起的角，$\e^{\i\varphi} = \cos\varphi + \i\sin\varphi$。
复共轭 $c\conj = u - \i v$ 使 $c\conj c = u^2 + v^2 = \abs{c}^2$。下一章将看到，沿选定的轴测量时，两个坐标的模平方就是两个出口的概率。
转动分析器后，新出口的概率由原来两分量的大小和相对相位共同决定。各章记号见卷首的\hyperref[chap:symbols]{常用符号表}。

全文采用 SI，$\hbar = h/(2\pi)$ 为约化 Planck 常数，$\varepsilon_0$ 为真空介电常数。
Bohr 磁子 $\muB = e\hbar/(2m_\electron)$ 给出电子磁矩的尺度，
精细结构常数 $\alpha = e^2/(4\pi\varepsilon_0\hbar c)$ 给出无量纲电磁耦合强度。

\reading{主要教材为Feynman第III卷与MIT 8.05、8.06讲义；书目及实验出处见书末。}
